\documentclass[11pt]{article}
\usepackage[utf8]{inputenc}
\usepackage[T1]{fontenc}
\usepackage{lmodern}
\usepackage[margin=1in]{geometry}
\usepackage{csquotes}
\usepackage[style=authoryear,backend=biber,sorting=nyt,maxbibnames=99,giveninits=true,uniquename=init]{biblatex}
\usepackage[hidelinks]{hyperref}
\addbibresource{references.bib}
\begin{document}
\section*{Drainage Management and Crop Yield Improvement}

\subsection{Overview}

Across a broad range of climates, soils, and crops, the evidence consistently establishes that agricultural subsurface drainage and its managed variants (controlled drainage, subirrigation, and drainage water recycling) raise crop yields relative to undrained or poorly drained conditions, primarily by relieving excess-water (waterlogging) stress early in the season and, where water tables can be managed, by alleviating water-deficiency stress later on. The strongest, most transferable conclusions come from long-term field experiments, multi-site and meta-analytic syntheses, and calibrated simulation studies. A United States meta-analysis of 154 pairwise comparisons found that controlled drainage and subirrigation increased crop yield by an average of \(8.0\%\) relative to conventional free drainage \parencite{Singh2022b}, while a combined experiment-and-producer study reported soybean yield gains of \(275\ \mathrm{kg\ ha^{-1}}\) (\(8\%\)) in experiments and \(157\ \mathrm{kg\ ha^{-1}}\) (\(4\%\)) in producer fields under artificial drainage \parencite{Mourtzinis2021}. The authors of that study framed subsurface drainage explicitly as an effective practice for increasing yields in environments prone to excess water early in the season \parencite{Mourtzinis2021}.

Several headline findings recur throughout the literature. First, yield benefits are typically larger in wet years and on poorly drained, fine-textured or naturally waterlogged soils, and they can be modest, absent, or even negative where existing drainage is already adequate or where over-drainage induces drought stress \parencite{JORJANI1991,Wiersma2010,Cordeiro2015,Kladivko2005}. Second, the response depends strongly on drain spacing, depth, and design, which mediate how quickly the water table is lowered \parencite{HARDJOAMIDJOJO1982,Skaggs2006,Rui2024,Acharya2019}. Third, in arid and semi-arid irrigated systems the yield benefit is often inseparable from salinity control and improved irrigation water-use efficiency \parencite{ELMOWELHI1988,Mathew2001,Yang2022,Liu2025}. Fourth, controlled drainage and subirrigation can retain water for the crop and deliver additional gains when precipitation is limiting \parencite{Kaur2021,Gunn2018,Gunn2019,Nelson2011}. This section synthesises these consolidated findings, the mechanisms behind them, and the quantified outcomes reported across the source studies.

\subsection{Magnitude of yield gains from subsurface drainage}

The clearest body of evidence concerns direct comparisons of drained versus undrained (or surface-drained/control) systems, which repeatedly document substantial yield increases. In a long-term Indiana experiment on Clermont soil, corn yields remained stable in the undrained control but progressively increased in all three drained treatments across 24 corn years, with the drainage effect becoming more pronounced in the last ten corn years than in the first ten; soybean yields, by contrast, were not significantly affected by drainage spacing \parencite{Rui2024}. A companion long-term Indiana dataset found whole-plot average corn yields of 9.8, 9.7, 9.5, and \(9.2\ \mathrm{Mg\ ha^{-1}}\) for the 5-, 10-, and 20-m spacings and the nondrained control, respectively, and concluded that drainage improvements are a long-term investment that may not yield benefits in every year \parencite{Kladivko2005}.

The largest relative gains are reported in heavily stressed environments. In the Nile Delta, subsurface drainage at 20-m spacing and 1.5-m depth doubled cotton and rice yields and increased wheat and clover yields by \(50\%\) over ten years relative to no subsurface drainage \parencite{ELMOWELHI1988}. In the acid saline "kari" soils of Kuttanad, India, subsurface drainage up to 30-m spacing raised average grain yield by \(1.36\ \mathrm{t\ ha^{-1}}\) (\(61\%\)) over the control, with grain yield and 100-grain weight significantly superior \parencite{Mathew2001}. In waterlogged saline clay soils of Bangladesh, drainage increased sunflower achene yield by \(30\textrm{--}64\%\) across drained treatments in two experiments, while in a third experiment the individual surface-drainage (SD) and subsoil-drainage (SSD) treatments increased yield by \(91\%\) and \(81\%\), respectively, and the combined subsoil-plus-surface drainage (SSD+SD) added only a further \(24\%\) and \(31\%\) increase relative to SD and SSD; an earlier companion study reported up to a \(95\%\) increase \parencite{Islam2025,Islam2022}. In Alberta, a shallow subsurface system on land out of production for 20 years raised barley grain yield from \(1310\ \mathrm{kg\ ha^{-1}}\) in 1977 to \(3900\ \mathrm{kg\ ha^{-1}}\) in 1978, with plant stand up \(43\%\) \parencite{BENNETT1982}.

More moderate but statistically robust gains dominate temperate rainfed systems. Soybean yield was \(8\%\) higher under artificial drainage in experiments and \(4\%\) higher in producer fields across the United States \parencite{Mourtzinis2021}, while shallow drainage in North Carolina increased corn yield by \(0.4\ \mathrm{Mg\ ha^{-1}}\) (\(4\%\), or \(6.6\%\) excluding an anomalous 2016 season) and soybean yield by \(0.46\ \mathrm{Mg\ ha^{-1}}\) (\(14.3\%\)) relative to free drainage \parencite{Watkins2024}. In Iranian subsurface-drained paddies, all drainage plots exceeded the control grain yield by \(13.4\textrm{--}31.3\%\) for rice \parencite{DarziNaftchali2017b}, and separate rice experiments reported grain yields \(22\textrm{--}66\%\) higher (2011) and \(32\textrm{--}70\%\) higher (2012) than the control, with the Bilevel treatment producing the highest yields \parencite{DarziNaftchali2014}. In China, blind-ditch and pipe drainage raised rice yield by \(8.0\%\) and \(6.0\%\) over the control, with the 3-m blind-ditch spacing producing the maximum yield of \(7.501\ \mathrm{Mg\ ha^{-1}}\) \parencite{Qi2025b}. Canola yield in Iran improved by roughly \(60\%\) under subsurface drainage relative to surface drainage, with mean grain yield higher by 697 and \(952\ \mathrm{kg\ ha^{-1}}\) in the two seasons \parencite{Asgari2021}. Cotton, sunflower, and potato experiments likewise reported increased actual yields under subsurface drainage \parencite{Yang2022,Zhang2024c,SerranoCoronel2018}.

\subsection{Controlled drainage, subirrigation, and drainage water recycling}

Managed drainage practices deliver yield benefits beyond free drainage when they retain water for crop use during dry periods. The U.S. meta-analysis quantified an \(8.0\%\) average yield increase from controlled drainage and subirrigation over conventional free drainage across 154 comparisons, drawn from 30 peer-reviewed articles and combined with fuzzy-logic mapping to identify suitable areas, including regions outside the well-studied Midwest \parencite{Singh2022b}. In Missouri, free drainage and drainage water recycling increased corn grain yield by \(20\%\) and \(26\%\), respectively, over a nondrained control averaged over 2015--2020, with recycling gains reaching \(17\%\), \(34\%\), and \(48\%\) in individual years and drainage water recycling positioned as a means to avoid drought losses in tile-drained landscapes \parencite{Kaur2021}. In North Carolina, drainage water recycling increased corn and soybean yields in all study years by \(0.52\ \mathrm{Mg\ ha^{-1}}\) (\(8\%\)) and \(0.45\ \mathrm{Mg\ ha^{-1}}\) (\(17\%\)); the soybean response was significant, subirrigation raised the groundwater table by an average of 15 cm, and the reservoir met irrigation demand in three of four seasons \parencite{Moursi2024}.

Subirrigation performance is strongly site-dependent, governed by lateral hydraulic conductivity. Across three Ohio sites, subirrigation raised the water table effectively at Fulton and Van Wert (growing-season water-table differences of 17--52 cm) and improved relative yield to \(87.3\textrm{--}93.2\%\) and \(80.4\textrm{--}90.8\%\), but had almost no hydrologic or yield effect at Defiance because of very low lateral hydraulic conductivity, where yields stayed near \(63\textrm{--}64\%\) \parencite{Gunn2018,Gunn2019}. Historic maize relative-yield responses to subirrigation ranged from essentially none at Defiance to \(26.5\%\) at Fulton, and simulations projected that subirrigation would increasingly outperform free drainage under future climate, raising its relative-yield advantage by \(2\%\), \(10\%\), and \(8\%\) at the three sites \parencite{Gunn2018}. In Minnesota, controlled drainage and subirrigation kept soil water content in optimal ranges, with the largest benefits observed in both a wet year (2013) and a dry year (2017), and subirrigation producing yields above the county average in a dry year \parencite{Niaghi2019}. In claypan-soil soybean systems, drainage plus subirrigation improved yield at both 6.1- and 12.2-m spacings with no difference between spacings, reduced interannual yield variability relative to drainage-only or nondrained controls, and allowed earlier planting where the nondrained control was delayed up to 17 days \parencite{Nelson2011}. Combining drainage water management with foliar pyraclostrobin produced synergistic soybean yield increases of up to \(1070\ \mathrm{kg\ ha^{-1}}\), with drainage-only or drainage-plus-subirrigation gains of \(18\textrm{--}22\%\) over the nondrained, nonfungicide control \parencite{Nelson2011b}. In Uzbekistan, controlled drainage managed the groundwater table so as to reduce water stress between irrigation events and increase grain yields, supporting the view that controlled drainage delivers good short-term returns through reduced water-deficiency stress \parencite{Dukhovny2018}.

\subsection{Drain spacing, depth, and system design}

Yield response is mediated by drainage intensity, set by spacing and depth. In the long-term Indiana experiment the narrowest 5-m spacing produced the highest average corn yield, followed by 10-m and 20-m spacings, with the undrained control lowest; narrower spacing also reduced the negative effect of post-planting rainfall on corn \parencite{Rui2024}. Distance-from-drain effects were more common at 20-m than at 10- or 5-m spacings, generally showing greater growth and yield near the tile \parencite{Kladivko2005}. On a Missouri claypan, the highest soybean yield occurred directly above the 6.1-m drainage tile line, all tile spacings and distances outyielded the nondrained control, and spatial yield variability was structured by distance from the tile line, suggesting narrower spacing to reduce variability in low-yielding environments \parencite{Nelson2017}.

Optimal design is soil- and rainfall-specific. Modelling for a Portsmouth sandy loam indicated that a maximum long-term relative corn yield of \(80\%\) required drain spacing of 40 m or less under good surface drainage but 30 m or less under poor surface drainage; at 100-m spacing, 25-year average relative yield fell to \(47\%\) (poor) versus \(61\%\) (good surface drainage) \parencite{HARDJOAMIDJOJO1982}. Across four soils, drain spacings that maximised returns ranged from 19 m (Normania) to 40 m (Portsmouth), with maximum returns near a drainage intensity of \(1.1\textrm{--}1.2\ \mathrm{cm\ day^{-1}}\) for most soils \parencite{Skaggs2006}. In Iowa, DRAINMOD predicted significant yield gains over undrained conditions for both Okoboji and Canisteo soils; Canisteo yields improved with closer spacing whereas Okoboji's largest improvement occurred at 18.3-m spacing \parencite{Arora2024}. In Minnesota, relative crop yield rose from \(43\%\) undrained to \(99\%\) at 23-m spacing and 135-cm depth, with profitability generally increasing with drain depth and spacing \parencite{Sands2015}. In the Canadian Prairies, however, over-drainage was a risk: a 15-m spacing could promote over-drainage and yield reductions, 40 m emerged as a threshold avoiding negative impacts on the sandy loam studied, and 60- and 90-m spacings were ineffective in wetter years \parencite{Cordeiro2015}. In North Dakota, spacing and depth effects depended on rainfall and crop: narrower (9 m) spacing gave the highest corn and soybean yields in the wet year and 3-year average, wider (15 m) spacing suited the dry year and sugarbeet, and a 1.2-m drain depth outperformed 0.9 m for sugarbeet \parencite{Acharya2019}. In China, deeper pipe burial and closer spacing generally raised yield, desalination, and irrigation water productivity, with 1.4-m burial identified as optimal in one region and D2 (45 cm) or a 45--60 cm range optimal under greenhouse conditions in another \parencite{Liu2025,Yang2022,Guo2025,Hou2018}. For sugarcane, a system able to maintain the water table between 0.3 and 1.5 m was concluded to increase cane and sugar yield \parencite{CAMP1982}.

\subsection{Mechanisms: waterlogging relief, aeration, roots, and physiology}

The proximate mechanism is relief of excess-water stress and restoration of root-zone aeration. Excess water quickly removed from the root zone provided sufficient aeration for barley growth \parencite{BENNETT1982}, and mid-season subsurface drainage lowered the water table in the rice root zone and improved aeration relative to the control \parencite{DarziNaftchali2014}. Growth-stage sensitivity matters: for cotton, the flowering and boll-forming stage was most waterlogging-sensitive and the budding stage most waterlogging-prone, and surface waterlogging caused greater yield loss than subsurface waterlogging, with combined stress less than additive \parencite{Qian2021}.

Improved aeration propagates through crop physiology and growth. Subsurface drainage increased cotton seedling emergence, root vigour, chlorophyll, plant height, leaf area, dry matter, and reproductive-organ weight, and path analysis identified seedling emergence rate as the largest direct contributor to yield—so drainage raised yield mainly by improving emergence through reduced soil salinity \parencite{Yang2022}. In sunflower, drainage decreased leaf \(\mathrm{Na^{+}}\) and \(\mathrm{Na^{+}/K^{+}}\) ratio and raised leaf \(\mathrm{K^{+}}\), stomatal conductance, and chlorophyll—all positively associated with achene yield—with early gains driven by waterlogging/salinity relief and late-season gains by improved water availability and solute potential \parencite{Islam2022}. In canola, subsurface drainage raised leaf area index, total dry weight, relative growth rate, and crop growth rate \parencite{Asgari2021}. In rice, subsurface drainage shifted biomass partitioning toward panicles, increased N uptake and translocation to reproductive tissue, and improved nutrient-use efficiency, raising NUE for grain by \(16.1\%\) and KUE by \(16\%\) \parencite{DarziNaftchali2022b}. Path analysis of blind-ditch rice systems attributed the largest direct contribution to yield to root biomass \parencite{Qi2025b}. A wheat root-trait review noted that dense adventitious roots, thin sparse roots, and aerenchyma-forming laterals are superior under waterlogging combined with other stresses, and that integrated subsurface tile-drained systems improve rhizosphere water status in poorly drained clay soils, with drainage among the agronomic strategies that relieve multiple-stress yield loss \parencite{Chen2025b}. Greenhouse and field work on root rot showed that watering intensity increased root rot and that no-drainage and shallow-drainage plots had greater root rot than conventional drainage in a wet year \parencite{Han2017}.

Drainage also raises crop water use where growth improves. In North Dakota, subsurface drainage increased peak growing-season evapotranspiration (e.g., 239 vs \(183\ \mathrm{mm}\) for July--August corn) and seasonal ET by \(16\%\) (2009) and \(7\%\) (2010), with higher crop coefficients in the drained field \parencite{Rijal2012}. Regionally, tile drainage slightly decreased annual cumulative ET (mainly early season) but drained areas showed slightly higher peak-season ET and generally higher LAI, indicating better crop growth \parencite{Yang2017,Yang2022}. Boundary-line analysis of cranberry fields identified an optimum water-table drawdown rate near \(1.25\ \mathrm{cm\ h^{-1}}\), with yield declining when drawdown deviated from that optimum, and soil hydrodynamic properties explaining nearly half the variance in drawdown and yield \parencite{Bulot2017}.

\subsection{Salinity control and reclamation in irrigated arid systems}

In arid and semi-arid irrigated agriculture, subsurface drainage improves yield largely by leaching salts and reclaiming soils. In the Nile Delta, one year of drainage reduced soil salinity in the top 1.5 m from 5.3 to \(2.2\ \mathrm{dS\ m^{-1}}\) without water beyond normal irrigation and rainfall \parencite{ELMOWELHI1988}. In Kuttanad, drainage leached sodium, calcium, and magnesium, drastically reduced chloride, lowered acidity, and homogenised soil chemistry within two years, alleviating the acidity, salinity, and cation/anion toxicity that had depressed rice yields \parencite{Mathew2001}. In Alberta, a shallow system controlled the water table (receding to pre-irrigation levels within 48 h) and reduced ECe and SAR mainly in the upper 0--0.3 m, removing \(2.2\ \mathrm{t\ ha^{-1}}\) of salt from the surface layer \parencite{BENNETT1982}. In Chinese saline systems, subsurface drainage improved root-zone desalination (\(13.54\%\) higher for the same irrigation), improved moisture uniformity, and raised sunflower and cotton yields, with deeper burial and closer spacing enhancing desalination \parencite{Zhang2024c,Yang2022,Liu2025,Guo2025}. Crop-straw drains (maize stalk, rice straw) removed salt and water more efficiently than perforated plastic pipes, improved desalination and root development, and improved maize grain yield—suggesting a low-cost organic substitute during initial reclamation of saline coastal soils \parencite{Lu2018}. Salinity management also has limits and risks: deficit irrigation in South Africa salinised upper soil layers, root-zone ECe exceeded the sugarcane threshold at all sites (implying \(>25\%\) yield reductions), and observed water-table depths and drainage discharges fell short of design, underscoring that installation alone is insufficient without maintenance \parencite{Malota2016}. Modelling for cotton showed that supplemental irrigation of \(\geq 100\ \mathrm{mm\ yr^{-1}}\) kept the root-zone salt balance favourable under a 1.2-m water-table control, whereas lower irrigation allowed salinity to accumulate \parencite{Li2018}. In Pakistan, water quality governed outcomes: good canal water maintained relative transpiration above 0.90 for wheat and cotton, whereas poor tubewell water dropped it as low as 0.60, with wheat more affected than cotton \parencite{Sarwar2001}.

\subsection{Water-use efficiency and irrigation water savings}

Controlled drainage repeatedly improves irrigation water-use efficiency and saves water without sacrificing yield. In Egypt, simple management concepts for existing subsurface systems, evaluated with DRAINMOD-S, improved irrigation water-use efficiency by \(15\textrm{--}20\%\) without yield reduction; crop-specific WUE gains were about \(16\%\) (wheat), \(23\%\) (maize), \(15\%\) (alfalfa), and \(20\%\) (cotton), with potential water savings of hundreds to over a thousand \(\mathrm{m^{3}\ ha^{-1}}\) per crop and roughly \(0.4\ \mathrm{BCM}\) over a two-year rotation in the Western Delta \parencite{Wahba2005}. A separate Egyptian study reported that controlled drainage cut applied irrigation water by \(15.6\%\) (wheat) and \(14\%\) (maize), raised water productivity to 1.03 and \(1.33\ \mathrm{kg\ m^{-3}}\), increased average maize relative yield by about \(6\%\), reduced drainage discharge by \(24.5\%\), and could save roughly \(1.33\ \mathrm{BCM}\) for El-Beheira governorate without affecting soil salinity \parencite{Mahmoud2021}. Modified drainage layouts reduced applied irrigation water by about \(30\textrm{--}40\%\) while restricted outflow had no adverse effect on soil salinity or crop yield for rice, maize, and cotton \parencite{ELATFY1991}. Controlled-water-table wheat in Egypt achieved the highest crop water productivity and large water savings at \(50\%\) available-water depletion with a 60-cm water table (identified as the best management practice) \parencite{ElHassanin2022}. In Iranian paddies, subsurface drainage raised water productivity and economic water productivity, with WUE increases of up to \(0.49\ \mathrm{kg\ m^{-3}}\) for one cultivar, while controlled subsurface systems increased water productivity without markedly harming irrigation application efficiency \parencite{DarziNaftchali2017b,Asgari2021}. In Pakistan, controlled subsurface drainage maintained groundwater near design depth, produced a maximum wheat yield of \(6.5\ \mathrm{t\ ha^{-1}}\) with water-use efficiency of \(0.93\ \mathrm{kg\ m^{-3}}\), and allowed large irrigation savings (23--129\% of the required depth) that could be reserved for drought \parencite{Khan2003}, while a water-conservation schedule could save up to \(25\%\) of scarce canal water \parencite{Sarwar2001}. In China, buried pipes raised irrigation water productivity over the control, highest at 1.4-m burial \parencite{Liu2025}.

\subsection{Weather dependence, timing, and earlier planting}

A consolidated finding is that drainage yield benefits scale with excess-water pressure and thus with weather. In Indiana, only the undrained control showed a significant negative corn-yield response to rainfall in the 14 days after planting, and growing-season rainfall was positively related to yield in drained systems, especially at 5-m spacing \parencite{Rui2024}. Wet years tended to favour narrower spacing while dry years favoured intermediate or wider spacings \parencite{Kladivko2005,Acharya2019}. Simulated yield reductions in the Canadian Prairies were more strongly associated with excess moisture early in the season than with late-season drought \parencite{Cordeiro2015}. For sugarcane, wet years favoured subsurface drainage whereas a dry year favoured a high, constant water table \parencite{CAMP1982}. In Japan, yield declined with the number of heavy-rain days (\(>30\ \mathrm{mm}\)) in fields without subsurface drainage but showed no such relationship where drainage was present, and mean buckwheat yield was much higher with subsurface drainage; raised-bed planting compensated for the absence of drainage except under exceptionally heavy rainfall \parencite{Takeshima2023}. Historically, comparisons on undrained versus artificially drained land showed corn yields more sensitive to precipitation on undrained land, with the largest physical yield gains from drainage on naturally poorly drained soils \parencite{JORJANI1991}.

Drainage also captures value through earlier field operations. Under producer conditions, artificial drainage was associated with sowing about 7--8 days earlier than natural drainage; combined with sowing-delay yield penalties of \(13\ \mathrm{kg\ ha^{-1}\ d^{-1}}\) (subsurface) and \(18\ \mathrm{kg\ ha^{-1}\ d^{-1}}\) (surface), one week's earlier sowing implied a \(91\textrm{--}126\ \mathrm{kg\ ha^{-1}}\) benefit \parencite{Mourtzinis2021}. Subsurface drainage similarly allowed earlier soybean planting and avoided the \(29\textrm{--}52\%\) harvest plant-population losses seen in nondrained controls in wet years \parencite{Nelson2011}.

\subsection{Modelling, design tools, and economics}

Simulation models—principally DRAINMOD—are well established for predicting drain flow, water-table depth, and relative yield, providing a validated basis for design. Reported performance is generally good: a modified DRAINMOD explained \(63\%\) of corn-yield variation over 12 years with close agreement on water-table elevations \parencite{HARDJOAMIDJOJO1982}; Indiana calibration/validation achieved Nash--Sutcliffe efficiencies of 0.70--0.85 for drain flow, with monthly flow predicted well in most plot-years and relative-yield patterns generally reproduced \parencite{Ale2009,Wang2006b}; Chinese predictions matched field water tables with a correlation of 0.97 \parencite{Li2018}; and Canadian Prairie simulations of water table and drainage outflow were rated good to excellent \parencite{Cordeiro2015}. In South Dakota, DRAINMOD reproduced measured subirrigated corn yields (representative \(R^{2}=0.53\)) and was judged the only assessment of DRAINMOD corn yields under predominantly irrigated conditions \parencite{DeBoer2006}. A simple steady-state Hooghoudt-based method gave reasonable estimates of drain depth and spacing to optimise profits, with a pooled drainage-rate regression (\(R^{2}=0.63\)) and soil-specific relationships up to \(R^{2}=0.94\) \parencite{Skaggs2006}. Detailed water-flow models (D-HYSAM) using a macroscopic root-extraction term reproduced soil-water pressure, moisture, and phreatic-surface positions in California \parencite{KARAJEH1994,KARAJEH1994b}. Design tools such as spreadsheet performance-index calculators enable multi-objective design that balances profitability against drainage volume and runoff \parencite{Sands2015}. Historically, the optimal water-table depth was recognised to depend on soil profile, crop, and permeability, and drain spacing could be indicated satisfactorily from soil permeability and root-zone conditions \parencite{HOOGHOUDT1952}.

Economically, drainage investments are frequently justified but soil- and crop-dependent. Iowa payback periods were seven years or less for all spacings, varying with soil type, spacing, and corn price \parencite{Arora2024}. Returns to drainage solely for higher corn yield were low on heavy soils (high installation cost) and highest on lighter soils, with the biggest payoff coming from enabling a shift from oats to corn on poorly drained land—so drainage is most attractive when it enables a switch from low- to high-value crops \parencite{JORJANI1991}. In Iran, installation costs were readily justified by annual increases in rice and canola yields, and subsurface drainage even enabled canola after rice (impossible under waterlogged surface drainage), raising both yield and cropping intensity \parencite{DarziNaftchali2014}. Combination subsurface irrigation-and-drainage systems in South Dakota appeared economically viable on suitable soils, with expected corn-yield increases of roughly \(40\textrm{--}100\%\) over dryland yields \parencite{DeBoer2006}. Managed strategies that minimise pumped water (e.g., continuous subirrigation from a defined growth stage with an optimised weir setting) reduced water use and labour costs while sustaining yield, enhancing the economic case for subirrigation \parencite{Ale2009,Gunn2019}.

\subsection{Agronomic side effects, integrated management, and adoption}

Beyond grain yield, drainage management affects grain quality, nutrients, and disease. Grain nutrient removal was greater under free drainage and drainage water recycling, though grain quality concentrations were largely unaffected in Missouri corn \parencite{Kaur2021}. In rice, milling recovery and protein content varied modestly across drainage treatments \parencite{DarziNaftchali2017b}. Notably, subsurface drainage that did not raise grain yield still increased grain protein: in spring wheat and soybean, protein rose with drainage coefficient even though yields were unchanged \parencite{Wiersma2010}. Foliar pyraclostrobin reduced Septoria brown spot and frogeye leaf-spot severity by \(2\textrm{--}8\%\) and interacted synergistically with drainage water management, while drainage itself did not affect disease severity \parencite{Nelson2011b}.

Drainage is increasingly embedded in integrated soil-and-water conservation. Reviews conclude that no single practice resolves difficult (e.g., duplex) soils, which require systems approaches combining tillage, rotations, ameliorants, and subsurface drainage; drainage was reported to eliminate waterlogging and deliver the largest, most reliable production gains on soils with stable subsoils \parencite{GARDNER1992}. Waterlogging-mitigation syntheses place combined drainage and crop management as the foremost step for severe waterlogging, with tolerant varieties and agronomic practices for minor cases, and the best strategy depending on severity, soil, crop, and local conditions \parencite{Manik2019}. Landscape-integrated frameworks combine tile drainage with terraces, water-retention structures, supplemental irrigation, and conservation tillage to improve field drainage while storing and reusing runoff and tile water \parencite{Li2021b}. Efficient drainage is also framed as a precondition for realising the benefits of conservation practices and improved genetics \parencite{Rui2024}, and cover crops or perennial forage can modestly reduce early-season drainage flow (e.g., a \(32\%\) May reduction under orchardgrass) with implications for flood mitigation in drained landscapes \parencite{Goeken2015,Daigh2014}. Broader land-use hydrology remains less certain, with catchment-scale flood impacts of drainage and land-use change identified as poorly understood and needing systems-thinking and connectivity-based upscaling \parencite{Rogger2017}.

Adoption evidence confirms drainage as a leading adaptive practice. Roughly three-fourths of surveyed U.S. producers reported at least some artificially drained cropland (22.6\% fully drained), and warm-season precipitation and experience with saturated soils were positive predictors of drainage adoption \parencite{Morton2015}. Finnish farmers prioritised well-functioning subsurface drainage and soil organic content as their main climate-adaptation measures—ahead of irrigation—with close to \(90\%\) rating subsurface drainage important, though many considered maintenance too expensive and called for subsidies \parencite{PeltonenSainio2021}. Japan's evolving rice-field drainage standards illustrate how design and policy co-evolve, including new uses of subsurface facilities for irrigation and ecosystem restoration \parencite{Ogino2007}.

\subsection{Boundary conditions and limits}

The evidence also delineates where drainage yields little or must be managed carefully. In the Red River of the North Basin, subsurface drainage did not improve spring wheat or soybean grain yield across drainage coefficients tested—only protein rose—because existing land leveling and surface drainage were already adequate, indicating limited applicability there \parencite{Wiersma2010}. Yield response also varies by soil texture: the meta-analysis found medium-textured soils most responsive to controlled drainage and subirrigation, with fine- and coarse-textured soils showing no clear effect, and the largest benefits under drier growing seasons in coarse and medium soils \parencite{Singh2022b}. Over-drainage can reduce yield through induced drought, as seen where 15-m spacing promoted yield loss and closer spacings and deeper weirs required careful management to avoid excessive water table lowering \parencite{Cordeiro2015,DeBoer2006}. Site hydraulic properties can nullify subirrigation benefits entirely, as at Defiance, Ohio \parencite{Gunn2018,Gunn2019}. Historic water-table experiments established the trade-off directly: yields increased as the water table was lowered from 40 to 150 cm, higher water levels produced earlier ripening and lower yield, and gains diminished at greater depths—so an optimum depth exists that varies with soil, crop, and permeability \parencite{HOOGHOUDT1952}. Finally, installation is necessary but not sufficient: without timely maintenance, observed water-table depths and drainage discharges can fall well short of design and fail to control shallow water tables \parencite{Malota2016}. Where soils are already well aerated—as in newly constructed sand rootzones—additional drainage produced no meaningful improvement in quality or rooting \parencite{Bigelow2001}. Collectively, these limits reinforce that the well-supported yield benefits of drainage are contingent on excess-water stress being present, on appropriate spacing and depth for the soil and climate, and on sustained system function.

\printbibliography
\end{document}
